\documentclass[11pt,a4paper,sans]{moderncv}
\moderncvstyle{classic}
\moderncvcolor{blue}
\usepackage[utf8]{inputenc}
% scale=0.84 (vs 0.8 in the resumes) keeps this letter to a single page
\usepackage[scale=0.84]{geometry}

\name{Alexandro}{Disla}
\title{Cover Letter}
\address{Rue Saint-Louis Jeanty \#14, Route de Frère}{Port-au-Prince, Haiti}{}
\phone[mobile]{(+509) 4148-3700}
\email{alexandrodisla@hotmail.com}
\extrainfo{GitHub: \url{https://github.com/AD0791}}

\begin{document}

\recipient{IOM Recruitment Team}{International Organization for Migration\\Vacancy CFA2026/023 -- Web Content Developer (Home-based)}
\date{\today}
\opening{Dear Hiring Manager,}
\closing{Sincerely,}
\enclosure[Attached]{Curriculum Vitae}

\makelettertitle

I am writing to apply for the position of \textbf{Web Content Developer} (CFA2026/023, International Consultant, Category B, home-based). I bring more than five years of fullstack development and digital content delivery experience, built almost entirely within remote consultancy arrangements and largely in the service of humanitarian and development programmes -- production web engineering, editorial discipline, and fluency in the operating context of an inter-governmental organization.

My technical foundation is directly relevant to web content work. I build and maintain sites and applications end to end: React and TypeScript on the front end, FastAPI and SQLAlchemy on the back end, with Docker and CI/CD pipelines for deployment. I develop and publish my own portfolio and technical writing platform using \textbf{Quarto}, a content-first publishing system, deployed to GitHub Pages with custom SCSS theming and light and dark variants. At Tekkod LLC I led a mobile modernization programme, upgrading a React Native codebase across major versions and migrating it to a new rendering architecture -- work demanding close attention to component structure, layout behaviour, and the user-facing quality of every screen.

Equally important for a content role, I understand that a website is a communication instrument before it is a codebase. Across my engagements with HANWASH, Caris Foundation International, and Anseye Pou Ayiti, my task was to take complex programmatic material -- indicator frameworks, survey results, monitoring data -- and render it legible for audiences with no interest in the underlying machinery: field teams, national partners, and donors. I produced technical guides, trained users, and built reporting interfaces in Power BI, Looker Studio, and Quarto that people actually used. Structuring information so that it is accurate, navigable, and appropriate to its reader is the same skill whether the output is a dashboard, a guide, or a web page.

I also work natively across the languages IOM operates in. I am a native speaker of \textbf{French} and \textbf{Haitian Creole} and fluent in \textbf{English}, and have written and reviewed professional content in all three -- so I can draft and adapt multilingual material directly rather than depending on a translation cycle.

Finally, I am accustomed to delivering as a home-based consultant. My work with HANWASH, Anseye Pou Ayiti, and Tekkod LLC has been fully remote, coordinated across time zones and delivered against defined outputs -- the model this Category B assignment assumes. My earlier engagement with CassionSoft on a \textbf{UNOPS} project gave me first-hand experience of United Nations procedures and documentation standards.

IOM's mandate depends heavily on how clearly its work is presented to the public, to member states, and to migrants themselves. I would welcome the opportunity to contribute to that effort, and would be glad to discuss how my background fits the requirements of this assignment. Thank you for your consideration.

\makeletterclosing

\end{document}
